\chapter{Opis technologiczny}

\section{API}

\subsection{Architektura aplikacji}

Backend aplikacji \verb|Wydatkomat| został zaimplementowany w architekturze REST API. System wykorzystuje model klient--serwer, w którym aplikacje webowa, mobilna oraz desktopowa komunikują się z centralnym backendem za pomocą protokołu HTTP.

Aplikacja została zaprojektowana zgodnie z architekturą warstwową. Logika systemu została podzielona na warstwy odpowiedzialne za:
\begin{itemize}
	\item obsługę żądań HTTP,
	\item logikę biznesową,
	\item komunikację z bazą danych,
	\item bezpieczeństwo i autoryzację użytkowników.
\end{itemize}

REST API udostępnia endpointy umożliwiające zarządzanie użytkownikami, wydatkami, płatnościami, znajomymi, powiadomieniami oraz statystykami. Dokumentacja endpointów została przygotowana zgodnie ze standardem OpenAPI i udostępniona przy użyciu Swagger UI.

\subsection{Technologie}

Backend aplikacji został napisany w języku Java w wersji 25 z wykorzystaniem frameworka Spring Boot 4. Do implementacji poszczególnych elementów systemu wykorzystano między innymi:
\begin{itemize}
	\item Spring Web -- obsługa REST API,
	\item Spring Security -- uwierzytelnianie i autoryzacja,
	\item Spring Data JPA -- komunikacja z bazą danych,
	\item Hibernate -- mapowanie obiektowo-relacyjne,
	\item PostgreSQL -- relacyjna baza danych,
	\item JWT -- obsługa tokenów autoryzacyjnych,
	\item OAuth2 -- logowanie zewnętrzne,
	\item Docker oraz Docker Compose -- konteneryzacja i wdrożenie aplikacji.
\end{itemize}

Do dokumentowania API wykorzystano standard OpenAPI oraz narzędzie Swagger UI.

\subsection{Baza danych}

System wykorzystuje relacyjną bazę danych PostgreSQL. Dane aplikacji przechowywane są w postaci encji powiązanych relacjami.

Komunikacja z bazą danych realizowana jest przy użyciu Spring Data JPA oraz Hibernate. Zastosowanie ORM pozwoliło ograniczyć liczbę bezpośrednich zapytań SQL oraz uprościć zarządzanie danymi.

Struktura bazy danych jest wersjonowana przy użyciu narzędzie \verb|Flyway|.

\subsection{Bezpieczeństwo}

Backend aplikacji został zabezpieczony przy użyciu mechanizmów dostępnych w Spring Security. System wykorzystuje uwierzytelnianie JWT, dzięki któremu możliwe jest bezstanowe zarządzanie sesją użytkownika.

Hasła użytkowników przechowywane są w postaci skrótów generowanych przy użyciu algorytmu Argon2. System wspiera również logowanie OAuth2 oraz uwierzytelnianie dwuetapowe z wykorzystaniem aplikacji Google Authenticator.

Komunikacja pomiędzy klientem a serwerem odbywa się wyłącznie przy użyciu protokołu HTTPS zabezpieczonego certyfikatem TLS 1.3. Dodatkowo system wykorzystuje mechanizmy ochrony przed CSRF oraz SQL Injection.

\subsection{Testy}

W projekcie zaimplementowano testy jednostkowe oraz integracyjne backendu. Testy wykorzystywane są do weryfikacji poprawności logiki biznesowej, autoryzacji użytkowników oraz działania endpointów REST API.

Do testów integracyjnych wykorzystano środowisko uruchamiane w kontenerach Docker przy użyciu narzędzia \verb|Testcontainers|. Pozwoliło to na testowanie aplikacji w środowisku zbliżonym do produkcyjnego.

Proces testowania obejmował między innymi:
\begin{itemize}
	\item poprawność uwierzytelniania użytkowników,
	\item działanie autoryzacji endpointów,
	\item walidację danych wejściowych,
	\item obsługę błędów,
	\item poprawność operacji na bazie danych.
\end{itemize}

\section{Aplikacja webowa}

Aplikacja webowa stanowi główny interfejs użytkownika systemu \verb|Wydatkomat|. Umożliwia zarządzanie wydatkami, znajomymi, płatnościami oraz profilem użytkownika bez konieczności instalowania dodatkowego oprogramowania.

Frontend został zaimplementowany przy użyciu biblioteki React z wykorzystaniem języka TypeScript oraz narzędzia Vite. Nawigacja pomiędzy widokami realizowana jest przy użyciu biblioteki React Router.

Komunikacja z backendem odbywa się przy użyciu REST API oraz protokołu HTTPS.

\section{Aplikacja mobilna}

Aplikacja mobilna została przygotowana dla systemu Android z wykorzystaniem frameworka React Native oraz środowiska Expo SDK 54. Interfejs aplikacji został zaimplementowany w języku TypeScript.

Do zarządzania nawigacją wykorzystano Expo Router, natomiast komunikacja z backendem realizowana jest przy użyciu biblioteki Axios. Dane uwierzytelniające użytkownika przechowywane są lokalnie przy użyciu Expo Secure Store.

Aplikacja mobilna umożliwia korzystanie z podstawowych funkcji systemu z poziomu urządzeń mobilnych.

\section{Aplikacja desktopowa}

Aplikacja desktopowa została napisana w języku Python i pełni głównie funkcję panelu administracyjnego systemu.

Klient desktopowy umożliwia zarządzanie użytkownikami systemu oraz wykonywanie operacji administracyjnych niedostępnych z poziomu aplikacji webowej i mobilnej.

Komunikacja aplikacji z backendem odbywa się przy użyciu REST API.